\section{Background and Related Work}\label{sec:background}

\paragraph{Chaos engineering.} Chaos engineering, as practised at Netflix, treats confidence in a distributed system as something to be gained by experiment: define a measurable steady state, hypothesise that it holds in both a control and an experimental group, introduce real-world events such as server failures into the experimental group, and look for a difference~\cite{basiri2016chaos}. The Principles of Chaos Engineering add that experiments should vary real-world events, run on production traffic where possible, be automated to run continuously, and keep the blast radius small~\cite{principlesofchaos}. Netflix's Chaos Automation Platform (ChAP) automated such experiments and compared experiment and control populations~\cite{blohowiak2016chap}. We borrow four ideas directly: the steady-state hypothesis, the paired control, fault injection on a real dependency path, and the aim of running experiments routinely. What changes for agents is the system under test: its steady state is a behaviour of a stochastic policy, not a throughput or error-rate metric.

\paragraph{Software fault injection.} Software fault injection deliberately introduces faults to assess how a system tolerates them, and has a long literature on fault models, injection points and representativeness~\cite{natella2016sfi}. Lineage-driven fault injection uses the lineage of successful outcomes to choose which faults to try~\cite{alvaro2015ldfi}, and Jepsen tests the safety claims of distributed databases under partitions and other faults~\cite{jepsen}. Our fault types (timeouts, errors, truncation, corruption) are classical; what is particular to agents is where they are injected (the agent's tool interface) and what is checked afterwards (the agent's next actions and its message to the user).

\paragraph{Tool-using agents and MCP.} ReAct interleaves reasoning and tool actions~\cite{yao2022react}; Toolformer trains a model to decide when to call APIs~\cite{schick2023toolformer}; Reflexion lets an agent revise its behaviour from verbal feedback~\cite{shinn2023reflexion}; and ToolLLM scales tool use to thousands of real-world APIs~\cite{qin2023toolllm}. MCP standardises how an agent client discovers and calls tools on a server, using JSON-RPC messages such as \texttt{tools/list} and \texttt{tools/call}~\cite{mcpspec}. Because MCP is a wire protocol with a small set of methods, a proxy that speaks it can sit between any conforming client and server without changing either. That property is what our method relies on.

\paragraph{Agent benchmarks and risk evaluation.} AgentBench~\cite{liu2023agentbench}, WebArena~\cite{zhou2023webarena}, SWE-bench~\cite{jimenez2023swebench}, $\tau$-bench~\cite{yao2024taubench} and ToolSandbox~\cite{lu2024toolsandbox} measure task success in environments that work as intended. $\tau$-bench also introduced the pass$^k$ metric, the probability that all of $k$ trials succeed, which makes run-to-run inconsistency visible~\cite{yao2024taubench}. ToolEmu uses a language model to emulate tools and another to judge the safety of an agent's actions~\cite{ruan2023toolemu}. InjecAgent~\cite{zhan2024injecagent} and AgentDojo~\cite{debenedetti2024agentdojo} evaluate indirect prompt injection, in which tool output carries an attacker's instructions. Our work differs from these in what it injects (operational faults rather than attacks) and in what it grades (behaviour after the fault rather than task success or attack success). Poisoned tool output is in our failure model, but only as planned work (\S\ref{sec:failure-model}).

\paragraph{Agents under tool faults.} Several 2026 studies address our question directly, and we position against them explicitly.
ReliabilityBench evaluates agents for consistency under repeated runs, robustness to task perturbations, and tolerance of tool faults (timeouts, rate limits, partial responses, schema drift), injected by a wrapper around simulated tools and graded by end-state equivalence~\cite{gupta2026reliabilitybench}.
ToolMisuseBench is an offline, deterministic benchmark with replayable faults and explicit step, call and retry budgets, reporting success, invalid calls, policy violations, recovery and efficiency~\cite{sigdel2026toolmisusebench}.
Kraishan wraps the environments of a function-calling benchmark in a fault-injection layer, compares against fault-free baselines, and finds that agents notice explicit errors far more often than plausible wrong values, and that they also report problems in a substantial share of fault-free trials~\cite{kraishan2026loud}.
AgentCheck records an agent's real MCP tool responses and replays the run with perturbed responses, grading with deterministic rules plus a model judge for interpretive labels~\cite{mazumder2026agentcheck}.
Zhang et al.\ inject ten fault types at tool, model, guardrail and inter-agent boundaries of multi-agent applications and release AgentChaosBench, a benchmark of telemetry traces with aligned no-fault controls, for detecting and localising faults after the fact~\cite{zhang2026agenticfail}.\footnote{Despite the similar names, AgentChaosBench~\cite{zhang2026agenticfail} and AgentChaos are unrelated. AgentChaosBench is a dataset for diagnosing faults from traces; AgentChaos is a tool for testing an agent's behaviour under injected faults.}
Closest to us, Iannillo et al.\ build MCP-Injector, proxy-style middleware on the MCP tool-call path between unmodified agents and real servers, with a fault model calibrated to incidents mined from MCP repositories and seedable fault schedules~\cite{iannillo2026mcpfaults}. Between the two gpt-oss backbones they test (20B and 120B parameters), they find no significant difference in fault tolerance. Their evaluation scores task completion with a rubric-based model judge. They also analyse behaviour: they flag runs with three or more consecutive calls to a failing tool, and they label ``tool bypass'', in which an agent gives a plausible answer without successful tool evidence, within a taxonomy of failure modes. Those labels are assigned after the fact from judge scores and traces, not checked against expectations declared before the run.

\paragraph{What is and is not new here.} Prior work already covers much of this space, and we are not the first to adapt chaos engineering to agents. ReliabilityBench describes its fault injection as chaos-engineering-style, grades deterministically by end-state equivalence with no model judge, and reports pass$^k$ over repeated runs~\cite{gupta2026reliabilitybench}. ToolMisuseBench is deterministic and budgets retries~\cite{sigdel2026toolmisusebench}. MCP-Injector puts a transparent proxy in front of real MCP servers and measures retry loops and tool bypass~\cite{iannillo2026mcpfaults}. AgentCheck's reproduce--intervene--mitigate loop is already a developer-testing workflow over MCP~\cite{mazumder2026agentcheck}. Kraishan~\cite{kraishan2026loud} and AgentChaosBench~\cite{zhang2026agenticfail} compare against fault-free runs. We claim none of these individually. Our claims are narrower. First, the expected safe behaviour is declared per experiment as deterministic predicates over behaviour (the call trace and the final answer) rather than over the end state, and it is checked on real tools behind a transparent proxy. Second, each experiment has a pass threshold over repeated runs and a paired no-fault control, so that the result can serve as a merge gate rather than a benchmark score. Third, the proxy is a reusable open core intended to host enforcement tools as well as fault injection. Whether this is practical, with verdicts stable enough and predicates precise enough to gate a merge, is an open empirical question, and the pilot is a first look at it.
